\chapter{Doctoral Entrance Problems}
\label{ch:problems}

The problems below are reproduced in French and routed by individual mathematical content.
\newpage

\section{Université Mohamed Khider --- Biskra}

\begin{problem}{(Biskra) — Épreuve N°02}{fa:biskra-2021-e2-p1}
Soit $H$ un espace de Hilbert et $F$ un sous-espace vectoriel fermé de $H$. On note $P_F:H\to H$ le projecteur orthogonal sur $F$, caractérisé par : pour tout $x\in H$, $P_F(x)\in F$ et $x-P_F(x)\in F^\perp$.
\begin{enumerate}\item Montrer que $F=\{x\in H : P_F(x)=x\}$.\item Montrer que $P_F$ est linéaire.\item Montrer que $P_F$ est continue et que $\|P_F\|=1$ (si $F\neq\{0\}$).\item Montrer que $\operatorname{Ker}(P_F)=F^\perp$.\end{enumerate}
\source{\emph{\small Université Mohamed Khider - Biskra, 2021, Épreuve~n\textsuperscript{o}~02.}}\end{problem}
\newpage

\begin{problem}{(Biskra) — Épreuve N°02}{fa:biskra-2021-e2-p2}
Soit $H=L^2([-1,1])$ muni du produit scalaire usuel $\langle f,g\rangle=\int_{-1}^1 f(t)g(t)\,dt$. On définit $T(f)=\int_{-1}^0f(t)\,dt-\int_0^1f(t)\,dt$.
\begin{enumerate}\item Montrer que $T$ est continue et calculer $\|T\|$.\item Déterminer $\operatorname{Ker}(T)$ et donner $g_0$ tel que $\operatorname{Ker}(T)=(\operatorname{Vect}\{g_0\})^\perp$.\item Calculer la projection orthogonale de $h(t)=t$ sur $\operatorname{Ker}(T)$.\end{enumerate}
\source{\emph{\small Université Mohamed Khider - Biskra, 2021, Épreuve~n\textsuperscript{o}~02.}}\end{problem}
\newpage

\begin{problem}{(Biskra) — Épreuve N°02}{fa:biskra-2021-e2-p3}
Soit $(X,\|\cdot\|_1)$ un espace de Banach et $(Y,\|\cdot\|_2)$ un espace normé. Soit $f:X\to Y$ continue et supposons qu'il existe $k>0$ tel que $k\|a-b\|_1\le\|f(a)-f(b)\|_2$ pour tous $a,b\in X$. Montrer que $f$ est une application fermée.
\source{\emph{\small Université Mohamed Khider - Biskra, 2021, Épreuve~n\textsuperscript{o}~02.}}\end{problem}
\newpage

\section{Université Ahmed Draïa --- Adrar}

\begin{problem}{(Adrar) — Épreuve N°03}{fa:adrar-2019-e3-p1}
On considère un problème aux limites non linéaire du second ordre sur $J=[0,b]$ avec conditions de Dirichlet et une non-linéarité $f(t,y)$. Montrer qu'il s'écrit sous la forme intégrale de Hammerstein
$$y(t)=\int_0^bG(t,s)f(s,y(s))\,ds,$$
où $G$ est la fonction de Green correspondante. Montrer que l'opérateur associé est continu et compact sur $E=C([0,b])$. Sous $|f(t,y)|\le\varphi(t)\psi(|y|)$, choisir une boule fermée convexe stable par l'opérateur et en déduire l'existence d'une solution par le théorème de Schauder.
\source{\emph{\small Université Ahmed Draïa - Adrar, 2019, Épreuve~n\textsuperscript{o}~03.}}\end{problem}
\newpage

\begin{problem}{(Adrar) — Épreuve N°03}{fa:adrar-2019-e3-p2}
Soit $E=C([0,1],\mathbb R)$ muni de $\|f\|_\infty=\sup_{x\in[0,1]}|f(x)|$. Montrer que $E$ est un espace de Banach. Pour $k\in C([0,1]^2)$, poser $\varphi_f(x)=\int_0^1k(x,t)f(t)\,dt$. Montrer que $\varphi_f\in E$ et que l'opérateur $K:f\mapsto\varphi_f$ est continu. Si $M=\sup_x\int_0^1|k(x,t)|\,dt<1$, montrer que
$$S=I-K+K^2-K^3+\cdots$$
converge dans $\mathcal L(E)$ et que $S=(I+K)^{-1}$.
\source{\emph{\small Université Ahmed Draïa - Adrar, 2019, Épreuve~n\textsuperscript{o}~03.}}\end{problem}
\newpage

\begin{problem}{(Adrar) — Épreuve N°03}{fa:adrar-2019-e3-p3}
Soit $I$ un intervalle ouvert de $\mathbb R$ et $1\le p<\infty$. Montrer qu'il existe une application linéaire continue $P:W^{1,p}(I)\to W^{1,p}(\mathbb R)$ telle que $(Pu)|_I=u$. Traiter les cas $I=]0,+\infty[$ et $I=]0,1[$ par réflexion et fonctions de coupure, puis conclure par changements de variables affines.
\source{\emph{\small Université Ahmed Draïa - Adrar, 2019, Épreuve~n\textsuperscript{o}~01.}}\end{problem}
